\scanpage{61}
% archive foliation: 60
% lines 1--8
\uline{Corollaire} Pour que $f$ soit un iso, il f. et s. que

a) $f$ induise
\[\pi_0(f):\pi_0(C)\overset{\sim}{\longrightarrow}\pi_0(C')\]
b) $f$ induise
\[\operatorname{Ob}f=f_0:\operatorname{Ob}(C)\simeq\operatorname{Ob}(\unsure{C'|C})\]
c) $\forall i\in I$, $f_{a_i,a_i}:\operatorname{Aut}(a_i)\simeq\operatorname{Aut}(f_0(a_i))$

% line 9
\uline{NB} $f$ équivalence $\Longleftrightarrow$ a) et c) sont vrais

% lines 10--14
\uline{Corollaire} Pour qu'une application
$f_0:C_0\longrightarrow C'_0$ provienne d'un
iso $f:C\simeq C'$, il f. et s. que

a) $f_0$ \deleted{\unsure{soit}} bijective

% lines 15--18
b) $f_0$ transforme la relation d'équivalence de $C_0$ (par comp. connexes)
en celle de $C'_0$ (condition vide si $C,C'$ connexes).

% lines 19--20
c) $\forall i\in I$, $\operatorname{Aut}(f_0(a_i))$ est isom. à $\operatorname{Aut}(a_i)$.

\medskip\centerline{\rule{2cm}{0.3pt}}\medskip
% lines 21--24
\deleted{\unsure{On}} Supposons $C=C'$. Alors on trouve
\[\operatorname{Aut}(C)\longrightarrow\unsure{\mathfrak{S}_{C_0}|\mathcal{S}_{C_0}}\qquad\uline{\text{surjective}}\]
si $\operatorname{Aut}^0(C)$ est le noyau, on trouve donc
